\clearpage
\section{Адаптивные интерфейсы и постановка исследовательской задачи}
\subsection{Понятия адаптивности, настраиваемости и персонализации}

Система, рассчитанная одновременно на начинающих и опытных пользователей, сталкивается с противоречием между простотой освоения и полнотой инструментов. Удаление сложных функций может упростить первые действия, но лишить специалистов необходимых возможностей. Отображение всех функций сразу создаёт обратный риск: повышает визуальную нагрузку и затрудняет поиск часто используемых операций. Разрешение этого противоречия связано не только с визуальным дизайном, но и с механизмами изменения состава и представления доступных функций.

В исследовании С.~Абрахао и соавторов проводится различие между адаптируемостью, когда пользователь самостоятельно изменяет интерфейс, и адаптивностью, когда изменение инициируется системой; также обсуждаются смешанные формы взаимодействия \cite{abrahao2021}. М.~Х.~Мираз, М.~Али и П.~С.~Экселл рассматривают адаптивные интерфейсы в контексте универсальной доступности и вариативности пользовательских потребностей \cite{miraz2021}. Эти положения позволяют отказаться от упрощённой трактовки адаптивности как автоматического скрытия кнопок: адаптация может касаться последовательности действий, терминологии, уровня детализации, представления данных и способа обратной связи.

Для SmartSQL целесообразно теоретически различать три режима. В \emph{начальном режиме} человек выполняет типовые действия с таблицей без знания SQL. В \emph{расширенном режиме} доступны фильтрация, соединения, работа со схемой и история операций. В \emph{профессиональном режиме} предусмотрены явный SQL и детальное управление запросами. Это предварительная проектная классификация, а не готовое решение и не результат тестирования пользователей. Важно, чтобы переключение режима не разрушало рабочий контекст и позволяло увидеть, какие именно действия будут выполнены.

\subsection{Модели пользователя и методы адаптации}

Л.~Ф.~Тагирова, А.~В.~Субботин и Т.~М.~Зубкова предлагают использовать характеристики пользователя и методы нечёткой логики для формирования персонализированного интерфейса \cite{tagirova2022}. Из этого следует практический для исследования вопрос: какие признаки действительно необходимы для выбора представления, а какие вводят избыточную сложность или нежелательную обработку персональной информации?

В рамках теоретической модели SmartSQL разумно начать с минимального набора признаков: выбранный пользователем уровень сложности, частота использования операций, контекст текущей задачи и явно выраженные предпочтения. Сбор иных личных характеристик не является обязательной предпосылкой адаптивности. Важно предусмотреть ручную отмену любых автоматических изменений, стабильность размещения критических действий и возможность вернуть интерфейс к известному состоянию.

Таким образом, автоматическая адаптация должна быть ограничена проверяемыми правилами. Например, система может предложить сократить набор второстепенных инструментов в режиме простого просмотра таблицы, но не должна скрывать предупреждение о массовом удалении данных. Наличие естественно-языкового ввода также не означает, что пользователь обязан постоянно общаться с системой через диалог: табличный и формальный способы работы следует сохранять как альтернативные каналы взаимодействия.

\subsection{Критерии качества и научная проверка адаптивности}

При оценке интерфейса недостаточно субъективно считать его «простым» или «красивым». А.~Н.~Морозов, В.~С.~Зарубин и С.~А.~Гришин исследуют количественные показатели информационного наполнения интерфейса и связь расположения элементов с профилем задачи \cite{morozov2020}. Ю.~Г.~Павлович и И.~Ф.~Киринович рассматривают A/B-тестирование как средство изменения пользовательского интерфейса с анализом поведения пользователей \cite{pavlovich2021}. Эти работы обосновывают необходимость эмпирической оценки, поскольку отдельное изменение оформления может оказывать противоположное влияние на разные группы задач.

Для последующих экспериментов предлагается сравнивать не абстрактные «старый» и «новый» интерфейсы, а воспроизводимые сценарии: создание таблицы, поиск записи, фильтрация по нескольким условиям, объединение данных, исправление ошибки и подтверждение изменения. Каждую задачу следует описывать одинаковой инструкцией и выполнять на сопоставимых наборах данных. Планируемые показатели включают время выполнения, долю успешно решённых задач, число ошибочных действий, количество обращений за подсказкой и субъективную оценку удобства.

\begin{table}[htbp]
\caption{Предлагаемые критерии будущей экспериментальной оценки}
\label{tab:metrics}
\small
\begin{tabularx}{\textwidth}{|p{42mm}|X|}
\hline
\textbf{Критерий} & \textbf{Способ наблюдения и интерпретации} \\\hline
Результативность & Доля задач, решённых правильно в пределах установленного сценария \\\hline
Временная эффективность & Время выполнения задачи с учётом уточнений и исправлений \\\hline
Ошибки пользователя & Количество неверных действий и их последствия \\\hline
Понимание операции & Способность объяснить, что произойдёт после подтверждения \\\hline
Стабильность навыка & Изменение времени и ошибок при повторном выполнении \\\hline
Удовлетворённость & Стандартизированный опрос после выполнения одинаковых задач \\\hline
\end{tabularx}
\end{table}

Представленные критерии в таблице~\ref{tab:metrics} являются программой оценки, а не результатами исследования. Для корректного эксперимента потребуется зафиксировать состав участников, уровень их опыта, порядок предъявления задач, условия обучения и способ обработки результатов. Особое внимание следует уделить эффекту порядка: интерфейс, который участник использовал вторым, может казаться удобнее только из-за предварительного знакомства с задачей.

\subsection{Синтез теоретических положений}

Научные публикации позволяют рассматривать адаптивность как управляемое изменение взаимодействия с учётом пользователя и задачи. Для SmartSQL это означает необходимость проектировать не один «универсальный экран», а согласованные сценарии, не скрывающие критическую информацию. Эффект адаптации должен подтверждаться экспериментально по измеримым показателям, а не только экспертной оценкой внешнего вида.